% --- 1. INFORMAZIONI GENERALI ---
\section{Informazioni Generali}
\begin{itemize}
    \item \textbf{Luogo/Piattaforma:} Server Discord\textsubscript{G}
    \item \textbf{Data:} 2026-06-09
    \item \textbf{Orario di inizio:} 14:30
    \item \textbf{Orario di fine:} 15:40
    \item \textbf{Responsabile\textsubscript{G}:} Francesco Marcon
    \item \textbf{Scriba:} Francesco Marcon
\end{itemize}

\vspace{0.5cm}
\textbf{Partecipanti presenti:}
\begin{table}[ht]
    \centering
    \renewcommand{\arraystretch}{1.2}
    \begin{tabularx}{\textwidth}{X l}
        \toprule
        \textbf{Nome e Cognome} & \textbf{Matricola} \\
        \midrule
        Sofia De Blasi & 2111014 \\
        Roberto Mariano Doroftei & 2111031 \\
        Filippo Idiometri & 2035617 \\
        Francesco Marcon & 2110990 \\
        Armando Moda Scarati & 2082864 \\
        Anton Ricardo Rupi & 2054304 \\
        \bottomrule
    \end{tabularx}
\end{table}

\vspace{0.5cm}
\textbf{Partecipanti assenti:}
\begin{table}[ht]
    \centering
    \renewcommand{\arraystretch}{1.2}
    \begin{tabularx}{\textwidth}{X l}
        \toprule
        \textbf{Nome e Cognome} & \textbf{Matricola} \\
        \midrule
        Nessuno & - \\
        \bottomrule
    \end{tabularx}
\end{table}

% --- 2. ORDINE DEL GIORNO ---
\section{Ordine del Giorno}
Gli argomenti previsti per la riunione odierna sono:
\begin{enumerate}
    \item Sprint\textsubscript{G} review e retrospective dello sprint 9,
    \item Controllo dello stato di avanzamento e compatibilità di frontend\textsubscript{G} e backend\textsubscript{G},
    \item Organizzazione delle presentazioni RTB\textsubscript{G} (generale e tecnologica) e pianificazione delle prove,
    \item Gestione dei rilasci del Proof of Concept\textsubscript{G} (PoC\textsubscript{G}), lettere di presentazione e aggiornamento sito web,
    \item Pianificazione dello sprint 10 ed allocazione delle task\textsubscript{G} e dei ruoli.
\end{enumerate}

% --- 3. RESOCONTO E DISCUSSIONE ---
\section{Resoconto della Riunione}

\subsection{Sprint review\textsubscript{G}, retrospective dello sprint 9 e allineamento documentale}
Durante la retrospective, il gruppo ha analizzato il ritardo di allineamento accumulato tra backend e frontend, causato dal tempo di attesa necessario per ricevere le chiavi API\textsubscript{G} da parte dell'azienda proponente. 
In fase di sprint review, si è constatato che il Piano di Progetto\textsubscript{G} (PdP) per lo sprint 9 è formalmente concluso. Per quanto riguarda il Piano di Qualifica\textsubscript{G} (PdQ), è stata evidenziata la necessità di eliminare alcuni refusi testuali; è stato inoltre stabilito che d'ora in avanti i grafici dovranno essere costantemente aggiornati e, qualora si presentino problematiche importanti, andranno integrate e dettagliate le relative descrizioni. 
Dal punto di vista dello sviluppo del codice, sia il frontend che il backend risultano funzionanti, evidenziando un unico problema di compatibilità specifico durante l'utilizzo del browser Brave. Infine, il gruppo ha convenuto sull'importanza di avviare un controllo globale e incrociato su tutta la documentazione, finalizzato alla correzione di errori ortografici, alla verifica della coerenza dei contenuti tra i diversi file e al completamento dei changelog.

\subsection{Pianificazione dello sprint 10 e preparazione RTB}
Il gruppo ha avviato la pianificazione dello Sprint 10. Una parte consistente della discussione è stata dedicata all'organizzazione della presentazione generale per l'RTB e alla preparazione della presentazione specifica sulle tecnologie. I membri del team hanno concordato di dividersi i rispettivi argomenti.
A livello operativo, sono state definite le attività di aggiornamento del sito web con l'inclusione del pacchetto compresso del Proof of Concept nel percorso \texttt{/poc/proof\_of\_concept.zip}. Parallelamente, verranno portate avanti la stesura della lettera di presentazione\textsubscript{G}, la compilazione del diario di bordo\textsubscript{G}, la rendicontazione finale delle ore effettive e l'avanzamento dei diagrammi e delle dashboard di PdP e PdQ per il nuovo sprint.
Per lo sprint 9, i ruoli e i compiti sono stati così distribuiti: Francesco Marcon opererà come Responsabile; Sofia De Blasi sarà impegnata come Progettista\textsubscript{G} sulle slide tecnologiche; Anton Ricardo Rupi e Roberto Mariano Doroftei gestiranno la documentazione e la presentazione RTB in qualità di Amministratori; Filippo Idiometri curerà la manutenzione del sito web come Amministratore\textsubscript{G}; infine, Armando Moda Scarati ricoprirà il ruolo\textsubscript{G} di Verificatore\textsubscript{G}.

% --- 4. DECISIONI PRESE ---
\section{Decisioni Prese}

\renewcommand{\arraystretch}{1.3}
\vspace{-0.3cm}
\noindent
\begin{longtable}{c p{12cm}}
    \toprule
    \textbf{Codice} & \textbf{Decisione} \\
    \midrule
    \endfirsthead

    \toprule
    \textbf{Codice} & \textbf{Decisione} \\
    \midrule
    \endhead

    \midrule
    \multicolumn{2}{r}{\textit{Continua nella pagina successiva}} \\
    \midrule
    \endfoot

    \bottomrule
    \endlastfoot

    \textbf{VI-17.1} & Approvazione del Piano di Progetto per lo sprint 9. \\[4pt]
    \textbf{VI-17.2} & Suddivisione delle tematiche per le presentazioni RTB, con successivo fissaggio di una data di prova collettiva. \\[4pt]
    \textbf{VI-17.3} & Approvazione del caricamento del file \texttt{proof\_of\_concept.zip} e contestuale aggiornamento del sito web. \\[4pt]
    \textbf{VI-17.4} & Redazione e finalizzazione della lettera di presentazione e del diario di bordo per la data del 12-06. \\[4pt]
\end{longtable}

% --- 5. AZIONI (TODO) ---
\section{Azioni da Svolgere (To-Do)}

\renewcommand{\arraystretch}{1.3}
\begin{longtable}{p{2.3cm} c p{4.5cm} p{3cm} l}
    \toprule
    \textbf{Codice} & \textbf{Rif. Decisione} & \textbf{Descrizione Task} & \textbf{Assegnatario} & \textbf{Scadenza} \\
    \midrule
    \endfirsthead

    \toprule
    \textbf{Codice} & \textbf{Rif. Decisione} & \textbf{Descrizione Task} & \textbf{Assegnatario} & \textbf{Scadenza} \\
    \midrule
    \endhead

    \midrule
    \multicolumn{5}{r}{\textit{Continua nella pagina successiva}} \\
    \midrule
    \endfoot

    \bottomrule
    \endlastfoot

    \ghissue{226} & VI-17.3 & Aggiornamento del sito web e inserimento del file \texttt{/poc/proof\_of\_concept.zip} & Filippo Idiometri & 2026-06-16 \\[4pt]

    \ghissue{219} & VI-17.2 & Creazione della presentazione RTB relativa alle tecnologie & Sofia De Blasi & 2026-06-16 \\[4pt]

    \ghissue{220} & VI-17.2 & Creazione della presentazione RTB generale & Roberto M. Doroftei & 2026-06-16 \\[4pt]

    \ghissue{221} & VI-17.4 & Stesura della lettera di presentazione del PoC & Francesco Marcon & 2026-06-16 \\[4pt]

    \ghissue{222} & VI-17.4 & Stesura del diario di bordo & Francesco Marcon & 2026-06-12 \\[4pt]

    \ghissue{225} & - & Stesura del verbale interno della riunione odierna & Francesco Marcon & 2026-06-16 \\[4pt]
\end{longtable}

% --- 6. PROSSIMA RIUNIONE ---
\section{Prossima Riunione}
\begin{itemize}
    \item \textbf{Data:} 2026-06-16
    \item \textbf{Ora:} 14:30
    \item \textbf{OdG preliminare\textsubscript{G}:} 
    \begin{itemize}
        \item Revisione\textsubscript{G} e verifica dei task del nono sprint;
        \item Simulazione e prove generali discorsive per le presentazioni RTB;
        \item Controllo finale della documentazione e approvazione della lettera di presentazione.
    \end{itemize}
\end{itemize}